本手册按《2027 年全国硕士研究生招生考试计算机学科专业基础考试大纲》的条目顺序编排，
覆盖大纲「考查内容」中数据结构、计算机组成原理、操作系统、计算机网络四科的全部条目。

\vspace{6pt}
{\bfseries 一、编排方式}

全书分四部分，共 24 章。四部分与大纲的四科一一对应，章与该科下的一级条目一一对应，
章内小节按大纲的二级、三级条目顺序展开。每一章的开头用引用块照抄本章对应的大纲原文条目，
便于对照检查覆盖情况；每一章末尾固定给出「本章小结表」与「易错清单」两张速查内容。

\vspace{6pt}
{\bfseries 二、阅读约定}

本手册用三种提示框把不同性质的内容区分开：

\begin{itemize}\setlength{\itemsep}{2pt}
\item \textbf{核心结论}——必须记住的定义、公式、判定条件，考前反复看的就是这些。
\item \textbf{易错点}——统考中反复出现的陷阱，以及容易与相邻概念混淆的地方。
\item \textbf{真题考法}——该知识点在统考中以什么形式出现；能确认年份与题号的，
      如实标注。无法确认的不标，不编造。
\end{itemize}

正文中的公式一律用数学排版；算法与代码用等宽框给出；对比性内容尽量做成表格。
凡涉及计算的知识点，都给出公式、符号含义、适用范围，并配一组可以自己动手复核的数字算例。

\vspace{6pt}
{\bfseries 三、素材来源}

本手册的素材来自三处：
\begin{itemize}\setlength{\itemsep}{2pt}
\item 语雀知识库「计算机考研 408 笔记」（月暗）导出的 182 篇笔记 Markdown 全文；
\item 工作区中的王道 2027 系列教材、课件、强化笔记与历年真题汇编；
\item 工作区 408 计算机考研目录下的各类复习笔记与答题本。
\end{itemize}

素材在整理过程中做了三件事：一是按大纲重新组织顺序；二是把原笔记中的口语化叙述、
错别字与图片占位改写成规范的知识点条目，公式按上下文重新推导书写；三是逐条核对，
凡素材之间口径冲突的，以大纲用词与统考真题为准，并在正文中说明差异。

\vspace{6pt}
{\bfseries 四、使用建议}

第一轮复习时，按章通读，重点看「核心结论」与例题，把「易错点」留到做题后再回看；
第二轮复习时只看每章的「本章小结表」，遇到想不起来的再翻回正文；
考前冲刺时只看「核心结论」与「易错清单」。

\vspace{10pt}
\noindent{\small 说明：本手册为复习资料汇编，用于知识点检索与自查，不替代教材与真题训练。}
